\documentclass[sigplan,screen]{acmart}

\settopmatter{printacmref=false}
\setcopyright{none}
\renewcommand\footnotetextcopyrightpermission[1]{}
\pagestyle{plain}

\usepackage{booktabs}
\usepackage{listings}
\usepackage{xcolor}
\lstset{
  basicstyle=\ttfamily\scriptsize,
  breaklines=true,
  frame=single
}

\AtBeginDocument{%
  \providecommand\BibTeX{{%
    \normalfont B\kern-0.5em{\scshape i\kern-0.25em b}\kern-0.8em\TeX}}}

\title{Team LongShot: Hybrid Retrieval with Neural Reranking for\\Temporally Evolving Scientific Collections \\ Progress Report}

\author{Adnan Harun Dogan}
\affiliation{%
  \institution{Middle East Technical University}
  \department{Department of Computer Engineering}
}
\email{adhd@ceng.metu.edu.tr}

\author{Umut Ozdemir}
\affiliation{%
  \institution{Middle East Technical University}
  \department{Department of Computer Engineering}
}
\email{umut@ceng.metu.edu.tr}

\begin{abstract}
We propose to participate in the CLEF 2026 LongEval-Sci shared task, which evaluates the temporal persistence of ad-hoc retrieval systems over evolving scientific document collections. Our approach progressively builds from classical course-based models (TF-IDF/VSM with SMART lnc.ltc scoring) through BM25 with query expansion, to dense semantic retrieval (Qwen3-Embedding) and neural cross-encoder reranking. We will investigate how each retrieval stage degrades across three temporal corpus snapshots and propose fusion techniques to mitigate temporal performance drops.
\end{abstract}

\maketitle

\section{Introduction}

Information retrieval systems are typically evaluated on static test collections, yet real-world document corpora are inherently dynamic: new documents are continuously published, existing ones are updated, and the information landscape shifts over time. This temporal evolution poses a fundamental challenge---a system that performs well on today's corpus may degrade as the collection grows and changes. The CLEF LongEval lab~\cite{longeval2023} was established precisely to study this phenomenon through longitudinal evaluation.

The LongEval-Sci task at CLEF 2026 operationalizes this challenge using the CORE open-access scientific literature collection~\cite{core2023}, which aggregates over 125 million research outputs from thousands of repositories worldwide. The task provides three temporal snapshots of this corpus (March--May, June--August, and September--November 2025) along with 338 test queries derived from real user search logs. Participants must submit ranked results for all three snapshots, and systems are evaluated not only on absolute retrieval quality but also on performance stability across time.


\section{Methodology/System Architecture}

Our system follows a four-stage progressive pipeline. Stages~1 (TF-IDF/VSM), 2 (BM25), and~3 (dense retrieval) are complete; Stage~4 (reranking) is upcoming. The overall data flow is shown below:

\begin{lstlisting}
Documents (JSONL)
  --> PyTerrier Indexer --> Inverted Index --> BM25 Retriever -->|
                                                                  |--> run.txt.gz
  --> Qwen3 Encoder -----> FlexIndex -------> Dense Retriever -->|

Queries (TSV) --> Tokenizer / Encoder --> Retriever
              --> TREC run file (per snapshot)
\end{lstlisting}

\textbf{Data.}
We use the abstract-only version of the LongEval-Sci 2026 corpus (585\,MB compressed, ${\approx}2$\,GB uncompressed). The corpus is distributed as 9 JSONL files containing 869,902 scholarly articles, each with title, abstract, author list, DOI, and publication metadata. Snapshot-1 is provided as training data; snapshots~2 and~3 are downloaded automatically via \texttt{ir-datasets-longeval} at runtime.

\textbf{Index.}
We use PyTerrier's \texttt{IterDictIndexer} backed by Terrier~5.11. Preprocessing follows Terrier defaults: whitespace tokenization, Porter stemming, standard English stopword removal, and case folding. We build a standard inverted index over the concatenation of title and abstract (\texttt{doc.default\_text()}). For the dense index, 4096-dimensional float32 embeddings are stored in a \texttt{FlexIndex} (pyterrier-dr library).

\textbf{Query Processing.}
Queries are single-field keyword queries loaded from TSV. For BM25, queries are tokenized identically to documents. For dense retrieval, queries are encoded with the Qwen3 instruction prefix (\texttt{prompt\_name="query"}) and documents with no prefix, following the model's asymmetric encoding design. Top-1000 documents are retrieved per query.

\textbf{User Interface.}
This is a shared task submission; there is no end-user interface. Output is standard TREC-format run files (\texttt{run.txt.gz}) per snapshot, to be submitted to the TIRA platform.

\subsection{Milestones and Schedule}

All planned work through Weeks~3--4 is now complete. Stage~1 (TF-IDF/VSM) was implemented and evaluated alongside BM25, and the dense retrieval run is underway. We are on track with the revised schedule.

\begin{table}[h]
\centering
\caption{Project schedule and current status.}
\label{tab:schedule}
\begin{tabular}{lll}
\toprule
\textbf{Week} & \textbf{Planned} & \textbf{Status} \\
\midrule
Weeks 1--2 (Apr 7--20)    & Environment setup, Stage 1 TF-IDF/VSM    & \textbf{Done} \\
Weeks 3--4 (Apr 21--May 3) & Stage 2 BM25, progress report           & \textbf{Done} \\
Weeks 5--6 (May 5--18)    & Stage 3 dense, Stage 4 reranking         & In progress  \\
Weeks 7--8 (May 19--Jun 2) & Temporal analysis, TIRA submission      & Upcoming     \\
Week 9 (Jun 2--7)          & Final report                             & Upcoming     \\
\bottomrule
\end{tabular}
\end{table}


\section{Experiments and Evaluation}

\subsection{Dataset}

All data has been acquired and prepared. The corpus consists of 869,902 open-access scientific articles from the CORE collection~\cite{core2023}.

\begin{table}[h]
\centering
\caption{Corpus and query statistics.}
\label{tab:corpus}
\begin{tabular}{lr}
\toprule
\textbf{Property} & \textbf{Value} \\
\midrule
Total documents (snapshot-1) & 869,902 \\
Document format              & JSONL (title + abstract) \\
Corpus size on disk          & ${\approx}2$\,GB (abstract-only) \\
Training queries             & 100 \\
Test queries                 & 338 (across 3 snapshots) \\
Training qrels --- raw       & 1,183 judgments (binary) \\
Training qrels --- DCTR      & 8,772 judgments (graded 0/1/2) \\
\bottomrule
\end{tabular}
\end{table}

Documents are loaded via \texttt{ir-datasets-longeval} using the dataset tag \texttt{longeval-sci-2026/clef-2026/sci}, which handles all three snapshots and their temporal metadata automatically.

\subsection{Inverted Index, TF-IDF/VSM, and BM25}

We have successfully built the inverted index and completed retrieval runs for both Stage~1 (TF-IDF/VSM) and Stage~2 (BM25) using PyTerrier~0.13. Run files have been produced for all three snapshots in valid TREC format.

\textbf{Toolkit.} PyTerrier~0.13 / Terrier~5.11, \texttt{ir-datasets-longeval}, SLURM cluster (64\,GB RAM, NVIDIA H200 node). BM25 completed in ${\approx}22$ minutes; TF-IDF/VSM in ${\approx}50$ minutes.

\textbf{TF-IDF/VSM (Stage~1).} We implement SMART lnc.ltc cosine scoring directly against the Terrier inverted index. Document vectors use log-tf with no idf and cosine normalisation (lnc); query vectors use log-tf, idf, and cosine normalisation (ltc). Porter stemming and English stopword removal are applied to queries to match the index preprocessing pipeline.

\textbf{Evaluation.} We evaluate on 100 training queries using \texttt{pytrec\_eval} with both raw (binary click) and DCTR (graded click-through rate) qrels. The primary metric is nDCG@10 as used by the shared task.

\begin{table}[h]
\centering
\caption{TF-IDF (lnc.ltc) vs.\ BM25 on 100 training queries (snapshot-1).}
\label{tab:results}
\begin{tabular}{llrrrr}
\toprule
\textbf{Qrels} & \textbf{Model} & \textbf{nDCG@10} & \textbf{nDCG@20} & \textbf{MAP} & \textbf{MRR} \\
\midrule
Raw  & TF-IDF (lnc.ltc) & 0.1856 & 0.2094 & 0.1494 & 0.3347 \\
Raw  & BM25             & 0.3310 & 0.3558 & 0.2658 & 0.4952 \\
\midrule
DCTR & TF-IDF (lnc.ltc) & 0.1638 & 0.1884 & 0.1426 & 0.2931 \\
DCTR & BM25             & 0.2922 & 0.3226 & 0.2573 & 0.4717 \\
\bottomrule
\end{tabular}
\end{table}

BM25 outperforms TF-IDF/VSM by a factor of ${\approx}1.8\times$ on nDCG@10. This gap is expected: BM25's document length normalisation via the $b$ parameter handles the variable-length scientific abstracts more effectively than cosine normalisation in lnc.ltc. The MRR of 0.495 for BM25 indicates a relevant document is typically found within the top~2 positions.

\textbf{Dense retrieval --- Stage~3 (in progress).}
We are currently running a dense retrieval baseline using Qwen3-Embedding-4B~\cite{qwen3} via \texttt{sentence-transformers} on GPU. Documents are encoded in batches of~32 on an NVIDIA H200 GPU. The model produces 4096-dimensional embeddings stored in a \texttt{FlexIndex} for exact nearest-neighbour retrieval (${\approx}3.5$ hours per snapshot). Snapshots~1 and~2 are complete; snapshot-3 is in progress. Upon completion we will evaluate and compare all three stages.

\subsection{Things to Do}

\begin{enumerate}
  \item \textbf{Stage~2 extensions:} Tune BM25 parameters ($k_1$, $b$) on training queries; implement RM3 query expansion.
  \item \textbf{Stage~3 --- Dense evaluation:} Complete the Qwen3-Embedding-4B run; compute nDCG@10 and compare with BM25 and TF-IDF.
  \item \textbf{Stage~3 --- Hybrid fusion:} Implement Reciprocal Rank Fusion (RRF) combining BM25 and dense scores.
  \item \textbf{Stage~4 --- Neural reranking:} Apply \texttt{ms-marco-MiniLM-L-12-v2} cross-encoder to top-100 candidates from the hybrid run.
  \item \textbf{Temporal stability analysis:} Measure nDCG@10 degradation across snapshots for each stage; compute Relative Performance Drop (RPD).
  \item \textbf{TIRA submission:} Submit all runs to the official evaluation platform.
  \item \textbf{Demo:} Build a simple query interface over the BM25 index.
\end{enumerate}

\bibliographystyle{ACM-Reference-Format}
\bibliography{references}

\end{document}
